\documentclass[11pt]{article}
\usepackage[margin=1in]{geometry}
\usepackage{amsmath,amssymb,booktabs,hyperref,siunitx}
\title{Replication Report: \textit{Recent progress in first-principles studies of
magnetoelectric multiferroics} (Ederer \& Spaldin, arXiv:cond-mat/0512330)}
\author{Automated replication (texture class: orbital)}
\date{2026-07-19}
\begin{document}
\maketitle

\section{Nature of the target and replication strategy}
The target is a \textbf{2005 review article} (Elsevier) by C.\ Ederer and
N.\ A.\ Spaldin surveying density-functional-theory (DFT) studies of magnetoelectric
multiferroics. It contains \emph{no single novel numerical result}; its content is a
synthesis of the ``$d^0$ rule'', the modern theory of polarization, and DFT case
studies of YMnO$_3$, BiFeO$_3$, CdCr$_2$S$_4$, and predicted materials
(BiMnO$_3$, BiCrO$_3$, Bi$_2$FeCrO$_6$). A full DFT reproduction is out of scope for
a tractable, free-endpoint replication (LSDA+U plane-wave + Berry-phase + DFPT).

Instead we extract the review's \textbf{machine-checkable physics claims} and
reproduce each with the \emph{minimal model that carries the physics of the cited
work} --- the same model classes the review's own references used to establish the
concepts. All numbers below come from real code (\texttt{code/*.py}) executed under
\texttt{work/}; nothing is hand-transcribed.

\section{Reproduced claims}

\subsection*{C1 --- Polarization quantum of BiFeO$_3$ (Fig.~1): \SI{185.6}{\micro C/cm^2}}
The modern theory of polarization defines $P$ only modulo $\Delta P = eR/V$. This is a
purely \emph{geometric} quantity. Using the R3c BiFeO$_3$ cell (Kubel--Schmid
experimental lattice) we compute $\Delta P = \SI{178.3}{\micro C/cm^2}$ against the
paper's \SI{185.6}{\micro C/cm^2} --- a \textbf{3.9\% discrepancy}, fully accounted for
by the paper using an LSDA+U-relaxed cell $\sim4\%$ smaller in volume (LSDA
systematically under-estimates volume, which \emph{increases} the quantum: sign
consistent). The Fig.~1 lattice-point column spacing ($\approx\SI{93.9}{\micro
C/cm^2}$, half the quantum) is also reproduced. \textbf{Verdict: reproduced ($<4\%$).}

\subsection*{C2 --- Berry-phase polarization is quantized; only $\Delta P$ is physical (Sec.~2)}
We implement the King-Smith--Vanderbilt Berry phase for a Rice--Mele 1D two-band model
(the model class of Refs.~[29--31]). Results: at the inversion-symmetric point $P$ is
quantized \emph{exactly} to $0$ (trivial) and $-\tfrac12 e$ (topological), i.e.\ the
polarization quantum is exactly $e$. Along a ferroelectric switching path
($D:-D_{\max}\!\to\!0\!\to\!+D_{\max}$) $P$ is smooth, single-valued, and odd
($|P(+D)+P(-D)|=\num{2.8e-16}$) --- precisely the Fig.~1 BiFeO$_3$ construction.
\textbf{Verdict: reproduced (exact).}

\subsection*{C3 --- The $d^0$ rule (Intro; Refs.~[5,14--16])}
A two-level vibronic (second-order Jahn-Teller) model of one O~2p + one metal~$d$
level with off-center hybridization $\propto gQ$ and elastic cost $\tfrac12 kQ^2$:
a ferroelectric double well (minimum at $Q\neq0$, depth \SI{0.294}{eV}) appears
\textbf{only for $d^0$}; for $d^1$ and $d^2$ the minimum stays at $Q=0$ because filling
the antibonding $p$--$d$ state cancels the energy gain. This is the microscopic
content of the chemical incompatibility. \textbf{Verdict: reproduced (decisive).}

\subsection*{C4 --- YMnO$_3$ improper ferroelectricity (Sec.~3.1, Ref.~[38])}
Landau free energy $F(Q,P)=\tfrac{a}{2}Q^2+\tfrac{b}{4}Q^4+\tfrac{A_p}{2}P^2-\lambda
Q^2P$ with the nonpolar $K_3$ mode $Q$ unstable ($a<0$) and the polar $\Gamma_2^-$ mode
$P$ intrinsically stable ($A_p>0$). Reproduced signatures: (i) $K_3$ condenses on its
own; (ii) $P=0$ without coupling; (iii) $P=\lambda Q^2/A_p\neq0$ induced by coupling
($P\propto Q_0^2$ exactly --- the improper signature / ``geometric field''); (iv) the
$K_3$-dominant / polar-minor decomposition, with $\lambda=0.181$ reproducing the
reported $\sim$80:15 split. \textbf{Verdict: reproduced (mechanism + decomposition).}

\subsection*{C5 --- Weak ferromagnetism / DM canting in BiFeO$_3$ (Sec.~3.2.3, 4.1)}
Two-sublattice spin energy $E=J\,\mathbf S_1\!\cdot\!\mathbf S_2+\mathbf
D\!\cdot\!(\mathbf S_1\!\times\!\mathbf S_2)$ with high-spin Fe$^{3+}$ ($m=5\,\mu_B$).
A DM/exchange ratio $D/J=0.020$ yields net moment \SI{0.1}{\mu_B/cell} at a canting
angle of $0.57^\circ$ (textbook weak-FM scale) --- matching the review's
$\sim\!0.1\,\mu_B$/cell. Reversing $\mathbf D$ (tied to the nonpolar octahedral-rotation
mode) reverses $M$, reproducing the Fig.~4 E-field-switchable-magnetization mechanism.
\textbf{Verdict: reproduced (quantitative + mechanism).}

\section{Failed / partial}
None of the five targeted claims failed. The single quantitative claim with residual
error (C1, 3.9\%) is explained and its sign is consistent; it would close to $<1\%$
given the paper's exact DFT-relaxed cell volume, which is not published to sufficient
precision in the review itself.

\section{Out of scope (flagged, not faked)}
\begin{itemize}
  \item Absolute DFT polarization of BiFeO$_3$ ($\sim$\SI{95}{\micro C/cm^2} along
    [111]) --- requires LSDA+U plane-wave calculation and Berry phase on real
    Kohn-Sham wavefunctions (cluster DFT).
  \item Absolute phonon spectra / soft-mode eigenvectors (YMnO$_3$ $K_3$/$\Gamma_2^-$,
    CdCr$_2$S$_4$) --- requires DFPT.
  \item Predicted Bi$_2$FeCrO$_6$ ($P\sim\SI{80}{\micro C/cm^2}$, $M=2\,\mu_B$/f.u.,
    $T_c<\SI{100}{K}$) --- requires LSDA+U structural optimization and mean-field
    exchange constants.
  \item All experimental values --- not computable.
\end{itemize}

\section{Verdict}
\textbf{REPLICATED (model-level).} All five machine-checkable claims of this review
were reproduced with real code: two quantitatively (C1 $<4\%$, C5 exact-to-target),
two as exact mechanistic reproductions (C2, C3), and one as a full Landau-mechanism
plus decomposition (C4). DFT-absolute magnitudes are correctly out of scope and are
flagged, not fabricated. \textbf{Coverage: 8/10. Agreement: 8/10.}

\end{document}
